into account the structure of @@M0@@ over @@M1@@, determined by 4.7, and the fact that the @@M2@@ are the classes of the sheaves occurring in the splitting of the inverse image of @@M3@@ on @@M4@@. This completes the proof of the first form of 4.6.

To deduce the second form, note that the elements @@M5@@ of the @@M6@@-@@M7@@-algebra @@M8@@ satisfy relations i) and ii). If @@M9@@ is the @@M10@@-@@M11@@-algebra generated by elements @@M12@@ satisfying i) and ii), there is therefore a @@M13@@-@@M14@@-homomorphism from @@M15@@ to @@M16@@ sending @@M17@@ to @@M18@@ (V 4.7). By 4.6, this homomorphism is surjective. Since the @@M19@@-ideal generated by relations i) and ii) in @@M20@@ contains the usual ideal generated by the relations of 4.6, this homomorphism is also injective, hence is an isomorphism.

\sgaStableTarget{sga6.e06.cor.4-8}\textbf{Corollary 4.8.} \uline{With the hypotheses and notation of} 4.5, \uline{put}:
@@M21@@
\uline{then the} @@M22@@-\uline{algebra} @@M23@@ \uline{is generated by the} @@M24@@; \uline{the ideal of relations among the} @@M25@@ \uline{is generated by the relations deduced from}:
@@M26@@

This is just another statement of 4.6. Indeed, the @@M27@@ can be expressed in terms of the @@M28@@, and conversely. Moreover, the relation: @@M29@@ is equivalent to the relation: @@M30@@.

\sgaStableTarget{sga6.e06.sec.4.9}4.9.
Let @@M31@@ be a scheme; denote by @@M32@@ the Grassmannian of type @@M33@@ of @@M34@@, with @@M35@@. As @@M36@@ varies, the @@M37@@ form an inductive system, the morphism @@M38@@ being defined by considering @@M39@@ as the
